\documentclass[12pt,a4paper]{article}

% --- Encodage et langue ---
\usepackage[utf8]{inputenc}
\usepackage[T1]{fontenc}
\usepackage[french]{babel}

% --- Police Times ---
\usepackage{mathptmx}

% --- Mise en page ---
\usepackage[top=2.5cm, bottom=2.5cm, left=2.5cm, right=2.5cm]{geometry}
\usepackage{setspace}
\onehalfspacing

% --- Graphiques et tableaux ---
\usepackage{graphicx}
\usepackage{float}
\usepackage{booktabs}
\usepackage{array}
\usepackage{multirow}
\usepackage{tabularx}
\usepackage{pgfplots}
\pgfplotsset{compat=1.18}
\usepackage{tikz}
\usetikzlibrary{shapes.geometric, arrows.meta, positioning, fit, backgrounds}
\usepackage{xurl}

% --- Maths et symboles ---
\usepackage{amsmath, amssymb}

% --- Liens ---
\usepackage[colorlinks=true, linkcolor=blue!60!black, citecolor=blue!60!black, urlcolor=blue!60!black]{hyperref}

% --- En-têtes ---
\usepackage{fancyhdr}
\pagestyle{fancy}
\fancyhf{}
\rhead{\textit{Projet NFP101 -- 2025/2026}}
\lhead{\textit{HandControlPC}}
\cfoot{\thepage}

% --- Légendes figures ---
\usepackage[font=small,labelfont=bf,textfont=it]{caption}

% --- Couleurs + boîtes ---
\usepackage{xcolor}
\usepackage{tcolorbox}
\tcbuselibrary{skins}

\definecolor{steelblue}{RGB}{70,130,180}
\definecolor{darkblue}{RGB}{25,60,120}

\pgfplotsset{
    rapportaxis/.style={
        tick label style={font=\small},
        label style={font=\small},
        title style={font=\small},
        grid=major,
        grid style={dashed, gray!30},
        scaled ticks=false,
    }
}

\newtcolorbox{descrbox}[1][]{%
  enhanced,
  colback=blue!3,
  colframe=blue!40!black,
  boxrule=0.5pt,
  arc=3pt,
  left=8pt, right=8pt, top=5pt, bottom=5pt,
  fontupper=\small,
  fonttitle=\small\bfseries\sffamily,
  title={\raisebox{-0.2ex}{\textsf{$\triangleright$}}~Point cl\'e},
  attach boxed title to top left={yshift=-2mm, xshift=4mm},
  boxed title style={colback=blue!40!black, colframe=blue!40!black, arc=2pt, boxrule=0pt},
  #1
}

\newtcolorbox{iabox}[1][]{%
  enhanced,
  colback=orange!5,
  colframe=orange!60!black,
  boxrule=0.5pt,
  arc=3pt,
  left=8pt, right=8pt, top=5pt, bottom=5pt,
  fontupper=\small,
  fonttitle=\small\bfseries\sffamily,
  title={\raisebox{-0.2ex}{\textsf{$\triangleright$}}~Utilisation IA},
  attach boxed title to top left={yshift=-2mm, xshift=4mm},
  boxed title style={colback=orange!60!black, colframe=orange!60!black, arc=2pt, boxrule=0pt},
  #1
}

% --- Enumération compacte ---
\usepackage{enumitem}
\setlist[enumerate]{topsep=2pt, itemsep=1pt, parsep=0pt}
\setlist[itemize]{topsep=2pt, itemsep=1pt, parsep=0pt}

% --- Listings code ---
\usepackage{listings}
\lstset{
  basicstyle=\ttfamily\small,
  backgroundcolor=\color{gray!8},
  frame=single,
  rulecolor=\color{gray!40},
  breaklines=true,
  showstringspaces=false,
  commentstyle=\color{gray},
  keywordstyle=\color{blue!70!black}\bfseries,
  stringstyle=\color{green!50!black},
  language=Python,
  numbers=left,
  numberstyle=\tiny\color{gray},
  xleftmargin=1.5em,
  framexleftmargin=1.5em,
}

% ============================================================
\begin{document}

% ============================================================
% PAGE DE TITRE
% ============================================================
\begin{titlepage}
\centering
\vspace*{1cm}
{\Large Conservatoire National des Arts et M\'etiers}\\[0.3cm]
{\large NFP101 --- Programmation Orient\'ee Objet}\\[3cm]

\rule{0.8\textwidth}{1.5pt}\\[0.6cm]
{\LARGE \textbf{HandControlPC}}\\[0.4cm]
{\Large Contr\^ole gestuel du PC par webcam}\\[0.2cm]
\rule{0.8\textwidth}{1.5pt}\\[2.5cm]

{\large Auteur : \textbf{Mohamed Amine Sobhi}}\\[0.4cm]
{\large Professeur : \textbf{Adrien Escourrou}}\\[2cm]

{\large Langage : \textbf{Python 3.11}}\\[0.3cm]
{\large Ann\'ee universitaire 2025--2026}\\[1.5cm]

\vfill
{\large Mai 2026}
\end{titlepage}

% ============================================================
\newpage
\tableofcontents
\newpage

% ============================================================
% RESUME
% ============================================================
\section*{R\'esum\'e}
\addcontentsline{toc}{section}{R\'esum\'e}

\noindent HandControlPC est une application Python de contr\^ole gestuel du PC par webcam.
Elle utilise MediaPipe pour d\'etecter 21 landmarks de la main en temps r\'eel et traduit
les gestes en actions syst\`eme macOS~: r\'eglage du volume, capture d'\'ecran, zoom vid\'eo,
et reconnaissance de signes (\'epeler \textbf{BONJOUR} lettre par lettre).

Le projet est structur\'e autour d'une hi\'erarchie de classes Python avec h\'eritage et polymorphisme,
une configuration externe JSON, 29 tests unitaires automatisables sans webcam,
et une documentation compl\`ete.

\vspace{0.3cm}
\noindent\textbf{Mots-cl\'es~:} Python, OpenCV, MediaPipe, programmation orient\'ee objet,
vision par ordinateur, interface gestuelle, tests unitaires.

\newpage

% ============================================================
% 1. CONTEXTE ET BESOIN
% ============================================================
\section{Contexte et besoin}
\label{sec:contexte}

\subsection{Probl\'ematique}

Les interfaces homme-machine traditionnelles (clavier, souris) imposent un contact physique
permanent avec le mat\'eriel. Certains sc\'enarios b\'en\'eficient d'une interaction sans contact~:
pr\'esentation en public, accessibilit\'e, environnements o\`u les mains sont occup\'ees.

HandControlPC propose une solution~: d\'etecter les gestes de la main via la webcam int\'egr\'ee
du PC et les convertir en actions syst\`eme, sans mat\'eriel suppl\'ementaire.

\subsection{Public cible}

\begin{itemize}
  \item Utilisateurs souhaitant contr\^oler le volume ou faire des captures sans quitter leur
        activit\'e principale.
  \item Personnes en situation de handicap moteur des membres sup\'erieurs (contr\^ole partiel).
  \item D\'eveloppeurs int\'eress\'es par la vision par ordinateur et les interfaces gestuelles.
\end{itemize}

\subsection{Cas d'usage principaux}

\begin{table}[H]
\centering
\caption{Fonctionnalit\'es principales et gestes associ\'es.}
\label{tab:usecases}
\small
\begin{tabularx}{\textwidth}{>{\bfseries}p{3.5cm} X p{4cm}}
\toprule
\textbf{Geste} & \textbf{Action syst\`eme} & \textbf{Retour au menu} \\
\midrule
Index seul lev\'e & Mode Volume~: pincement pouce-index ajuste le volume (0--100\,\%) & Main compl\`etement ouverte \\
\midrule
3 doigts lev\'es & Mode Screenshot~: maintenir le geste $\sim$0,7\,s d\'eclenche une capture & Automatique apr\`es capture \\
\midrule
2 mains visibles & Mode Zoom~: \'ecartement/rapprochement des index zoome le flux vid\'eo & Une seule main restante \\
\midrule
S\'equence B-O-N-J-O-U-R & Reconnaissance de signes~: \'epeler BONJOUR lettre par lettre & Touche \texttt{R} / automatique \\
\bottomrule
\end{tabularx}
\end{table}

\newpage

% ============================================================
% 2. ARCHITECTURE ET CHOIX TECHNIQUES
% ============================================================
\section{Architecture et choix techniques}
\label{sec:architecture}

\subsection{Structure des fichiers}

\begin{table}[H]
\centering
\caption{Organisation des modules du projet.}
\label{tab:modules}
\small
\begin{tabularx}{\textwidth}{>{\ttfamily}p{5.5cm} X}
\toprule
\textbf{Fichier} & \textbf{R\^ole} \\
\midrule
main.py & Point d'entr\'ee --- boucle principale, gestion des modes \\
test\_sign\_recognition.py & D\'emo reconnaissance de signes et s\'equence BONJOUR \\
config.json & Configuration externe (cam\'era, seuils, dossiers) \\
requirements.txt & D\'ependances Python \\
\midrule
utils/gesture\_action.py & Classe abstraite \texttt{GestureAction} (ABC) \\
utils/hand\_tracking.py & \texttt{HandDetector} --- d\'etection MediaPipe \\
utils/volume\_control.py & \texttt{VolumeController} --- h\'erite de \texttt{GestureAction} \\
utils/screenshot.py & \texttt{ScreenshotTaker} --- h\'erite de \texttt{GestureAction} \\
utils/zoom\_control.py & \texttt{ZoomManager} --- h\'erite de \texttt{GestureAction} \\
utils/config\_loader.py & Chargement JSON avec deepcopy et fallback \\
\midrule
tests/test\_gesture\_logic.py & 19 tests logique gestuelle \\
tests/test\_config.py & 7 tests configuration \\
tests/test\_screenshot.py & 3 tests h\'eritage et dossiers \\
\midrule
docs/ & Documentation PDF et LaTeX \\
screenshots/ & Captures d'\'ecran g\'en\'er\'ees \\
\bottomrule
\end{tabularx}
\end{table}

\subsection{Hi\'erarchie de classes et polymorphisme}

Le projet repose sur une classe abstraite \texttt{GestureAction} dont h\'eritent les trois
actions gestuelles. Cette conception respecte le principe \textbf{Open/Closed}~: ajouter
une nouvelle action ne n\'ecessite que de cr\'eer une nouvelle sous-classe.

\begin{figure}[H]
\centering
\begin{tikzpicture}[
  node distance=1.4cm and 2.2cm,
  classbox/.style={
    rectangle, draw=darkblue, fill=blue!6,
    rounded corners=3pt, minimum width=3.8cm, minimum height=1cm,
    align=center, font=\small\ttfamily
  },
  abstract/.style={
    classbox, fill=steelblue!20, draw=steelblue,
    font=\small\ttfamily\bfseries
  },
  arrow/.style={-{Latex[length=3mm]}, thick, draw=darkblue}
]

\node[abstract] (ga) {GestureAction\\{\tiny\normalfont (ABC)}};

\node[classbox, below left=of ga]  (vc) {VolumeController};
\node[classbox, below=of ga]       (st) {ScreenshotTaker};
\node[classbox, below right=of ga] (zm) {ZoomManager};

\draw[arrow] (vc) -- (ga);
\draw[arrow] (st) -- (ga);
\draw[arrow] (zm) -- (ga);

% Méthodes GestureAction
\node[right=3.5cm of ga, align=left, font=\small] (methods) {
  \textit{M\'ethodes abstraites~:}\\
  \quad \texttt{+ process(img, hands)}\\[4pt]
  \textit{M\'ethodes concr\`etes~:}\\
  \quad \texttt{+ activate()}\\
  \quad \texttt{+ deactivate()}\\
  \quad \texttt{+ render\_hud(img, text)}
};
\draw[dashed, gray] (ga.east) -- (methods.west);

\end{tikzpicture}
\caption{Hi\'erarchie d'h\'eritage~: \texttt{GestureAction} est la classe abstraite commune.
Chaque sous-classe impl\'emente \texttt{process()} selon sa logique propre (polymorphisme).}
\label{fig:hierarchy}
\end{figure}

\begin{descrbox}
Le \textbf{polymorphisme} est exploit\'e dans \texttt{main.py}~: la boucle principale peut
th\'eoriquement appeler \texttt{action.process(img, hands\_data)} sans conna\^itre le type
concret de l'action. La m\'ethode \texttt{render\_hud()} partag\'ee affiche automatiquement
le nom du mode actif en overlay sur le flux vid\'eo.
\end{descrbox}

\subsection{Pipeline de traitement vid\'eo}

\begin{figure}[H]
\centering
\begin{tikzpicture}[
  node distance=0.4cm,
  step/.style={rectangle, draw=darkblue, fill=blue!6, rounded corners=3pt,
               minimum width=2.8cm, minimum height=0.8cm, align=center, font=\small},
  arrow/.style={-{Latex[length=2.5mm]}, thick, draw=darkblue}
]
\node[step] (cam)   {Webcam\\(OpenCV)};
\node[step, right=of cam]  (mp)    {D\'etection\\(MediaPipe)};
\node[step, right=of mp]   (lm)    {21 landmarks\\$(x,y)$};
\node[step, right=of lm]   (fp)    {\texttt{fingers\_up()}};
\node[step, right=of fp]   (mode)  {Mode actif};
\node[step, right=of mode] (sys)   {Action\\syst\`eme};

\draw[arrow] (cam)  -- (mp);
\draw[arrow] (mp)   -- (lm);
\draw[arrow] (lm)   -- (fp);
\draw[arrow] (fp)   -- (mode);
\draw[arrow] (mode) -- (sys);
\end{tikzpicture}
\caption{Pipeline de traitement~: chaque frame passe de la webcam \`a une action syst\`eme
en moins de 33\,ms (\textgreater\,30\,FPS).}
\label{fig:pipeline}
\end{figure}

\subsection{Syst\`eme de configuration}

Tous les param\`etres sont externalis\'es dans \texttt{config.json}, charg\'e au
d\'emarrage par \texttt{config\_loader.py}. Un \texttt{deepcopy} prot\`ege le dictionnaire
par d\'efaut contre les modifications en place (bug corrig\'e lors du d\'eveloppement).

\begin{lstlisting}[caption={Extrait de config.json}, label={lst:config}]
{
  "camera":    { "width": 1280, "height": 720, "index": 0 },
  "detection": { "min_detection_confidence": 0.7,
                 "gesture_history_length": 15 },
  "volume":    { "min": 0, "max": 100,
                 "distance_min": 50, "distance_max": 300 },
  "screenshot":{ "folder": "screenshots", "stable_threshold": 0.6 },
  "zoom":      { "min_scale": 0.5, "max_scale": 3.0 }
}
\end{lstlisting}

\subsection{Biblioth\`eques utilis\'ees}

\begin{table}[H]
\centering
\caption{Biblioth\`eques et justification du choix.}
\label{tab:libs}
\small
\begin{tabularx}{\textwidth}{>{\bfseries}p{3cm} X}
\toprule
\textbf{Biblioth\`eque} & \textbf{R\^ole et justification} \\
\midrule
OpenCV & Capture webcam, traitement d'image, affichage HUD en temps r\'eel. Standard industriel,
        tr\`es performant sur macOS. \\
MediaPipe & D\'etection de 21 landmarks de la main par r\'eseau de neurones (Google).
           Fonctionne sur CPU standard, 30+ FPS. \\
NumPy & Interpolation lin\'eaire (\texttt{np.interp} pour le volume), clipping du zoom
       (\texttt{np.clip}). \\
MSS & Capture d'\'ecran multi-plateforme, plus rapide que PyAutoGUI pour ce cas. \\
osascript & API macOS native pour le contr\^ole du volume et le retour sonore. \\
pytest & Framework de tests unitaires~: 29 tests automatisables sans webcam. \\
\bottomrule
\end{tabularx}
\end{table}

\newpage

% ============================================================
% 3. FONCTIONNALITES DETAILLEES
% ============================================================
\section{Fonctionnalit\'es d\'etaill\'ees}
\label{sec:fonctionnalites}

\subsection{Mode Volume}

Le volume est r\'egl\'e proportionnellement \`a la distance entre le pouce (landmark~4)
et l'index (landmark~8)~:

\begin{equation}
  V = \left\lfloor \text{interp}\!\left(d_{\text{pouce-index}},\;
      [d_{\min},\, d_{\max}],\; [0,\, 100]\right) \right\rfloor
  \label{eq:volume}
\end{equation}

avec $d_{\min} = 50$ px et $d_{\max} = 300$ px (param\'etrables dans \texttt{config.json}).
L'appel syst\`eme \texttt{osascript} applique le volume~; une ligne verte relie les deux
doigts comme retour visuel. \textbf{Retour au menu~:} main compl\`etement ouverte (5 doigts).

\subsection{Mode Screenshot}

Un historique glissant de \texttt{GESTURE\_HISTORY\_LENGTH} frames est maintenu.
La capture n'est d\'eclench\'ee que si le taux de frames avec 3 doigts lev\'es
d\'epasse le seuil \texttt{stable\_threshold}~=~0,6~:

\begin{equation}
  \text{confiance} = \frac{\sum_{i=1}^{N} \mathbf{1}[\text{geste valide}_i]}{N}
  \geq 0{,}6
  \label{eq:confidence}
\end{equation}

Ce m\'ecanisme \'evite les faux positifs lors des transitions. Un son de confirmation
(\texttt{afplay}) et un message \texttt{CAPTURE REUSSIE!} confirment visuellement l'action.

\subsection{Mode Zoom}

Le facteur de zoom est calcul\'e \`a partir de la variation de distance entre les index
des deux mains~:

\begin{equation}
  \text{scale} = \text{clip}\!\left(\frac{d_\text{actuelle}}{d_\text{initiale}},\;
      s_{\min},\, s_{\max}\right)
  \label{eq:zoom}
\end{equation}

Une transformation affine \texttt{cv2.warpAffine} applique le zoom centr\'e sur le
milieu des deux mains.

\subsection{Reconnaissance de signes --- S\'equence BONJOUR}

\begin{table}[H]
\centering
\caption{Patterns de doigts pour les 8 signes reconnus~:
[pouce, index, majeur, annulaire, auriculaire].}
\label{tab:signs}
\small
\begin{tabular}{clp{6cm}}
\toprule
\textbf{Lettre} & \textbf{Pattern binaire} & \textbf{Geste physique} \\
\midrule
A & \texttt{[1, 0, 0, 0, 0]} & Poing ferm\'e, pouce sur le c\^ot\'e \\
B & \texttt{[0, 1, 1, 1, 1]} & 4 doigts lev\'es, pouce repli\'e \\
J & \texttt{[1, 0, 0, 0, 1]} & Pouce + auriculaire (signe \textit{shaka}) \\
L & \texttt{[1, 1, 0, 0, 0]} & Pouce + index (pistolet) \\
N & \texttt{[1, 1, 1, 0, 0]} & Pouce + index + majeur \\
O & \texttt{[0, 1, 1, 1, 0]} & Index + majeur + annulaire (3 doigts du milieu) \\
R & \texttt{[0, 1, 0, 0, 0]} & Index seul point\'e \\
U & \texttt{[0, 1, 1, 0, 0]} & Index + majeur (signe V / paix) \\
\bottomrule
\end{tabular}
\end{table}

Le programme attend que l'utilisateur \'epelle \textbf{B\,--\,O\,--\,N\,--\,J\,--\,O\,--\,U\,--\,R}
dans l'ordre. Chaque lettre valide apr\`es \texttt{HOLD\_FRAMES}~=~20 frames stables.
Une erreur de s\'equence remet le compteur \`a z\'ero. Quand la s\'equence est compl\`ete,
le programme \'emet vocalement \og{}Bonjour\fg{} et affiche un message de succ\`es.

\begin{descrbox}
Le choix des patterns a \'et\'e it\'eratif~: les premiers patterns (auriculaire seul,
pouce+majeur sans index) \'etaient anatomiquement difficiles \`a tenir et peu fiables.
Les patterns actuels ont \'et\'e s\'electionn\'es pour \^etre \textbf{distincts} entre eux
et \textbf{naturels} \`a ex\'ecuter, permettant une d\'etection stable.
\end{descrbox}

\newpage

% ============================================================
% 4. TESTS UNITAIRES
% ============================================================
\section{Validation et tests}
\label{sec:tests}

\subsection{Strat\'egie de test}

Tous les tests sont con\c{c}us pour s'ex\'ecuter \textbf{sans webcam, sans mat\'eriel, sans appel
syst\`eme r\'eel}. Les d\'ependances externes (MediaPipe, MSS, osascript) sont soit
mock\'ees avec \texttt{unittest.mock}, soit contournables par injection de donn\'ees.

\subsection{Couverture}

\begin{table}[H]
\centering
\caption{29 tests unitaires r\'epartis en 3 fichiers.}
\label{tab:tests}
\small
\begin{tabularx}{\textwidth}{>{\ttfamily}p{5cm} c X}
\toprule
\textbf{Fichier} & \textbf{Tests} & \textbf{Ce qui est test\'e} \\
\midrule
test\_gesture\_logic.py & 19 & Reconnaissance de signes A, B, J, L, N, O, R, U~;
                               s\'equence BONJOUR~; calcul de volume~;
                               ZoomManager (reset, clipping)~; fingers\_up \\
test\_config.py         &  7 & Chargement JSON valide, JSON invalide, fichier manquant,
                               d\'efauts, deepcopy, cl\'es pr\'esentes \\
test\_screenshot.py     &  3 & Cr\'eation dossier, h\'eritage GestureAction, mock MSS \\
\midrule
\textbf{Total}          & \textbf{29} & \textbf{100\,\% passent sans webcam} \\
\bottomrule
\end{tabularx}
\end{table}

\begin{figure}[H]
\centering
\begin{tikzpicture}
\begin{axis}[
    rapportaxis,
    ybar,
    width=10cm, height=5.5cm,
    symbolic x coords={test\_gesture\_logic, test\_config, test\_screenshot},
    xtick=data,
    xticklabel style={font=\small, rotate=10},
    ylabel={Nombre de tests},
    title={R\'epartition des tests par fichier (29 tests --- 100\,\% PASSED)},
    nodes near coords,
    nodes near coords align={vertical},
    every node near coord/.append style={font=\small, yshift=2pt},
    ymin=0, ymax=25,
    bar width=1.5cm,
    fill=steelblue!60,
    draw=steelblue!80,
]
\addplot coordinates {
    (test\_gesture\_logic, 19)
    (test\_config, 7)
    (test\_screenshot, 3)
};
\end{axis}
\end{tikzpicture}
\caption{29 tests unitaires r\'epartis en 3 modules. Ex\'ecution~: \texttt{python -m pytest tests/ -v}.}
\label{fig:tests}
\end{figure}

\subsection{Exemple de test unitaire}

\begin{lstlisting}[caption={Test de la reconnaissance de signe et de la s\'equence BONJOUR.},
                   label={lst:test}]
def test_recognize_sign_B():
    assert recognize_sign([0, 1, 1, 1, 1]) == "B"

def test_bonjour_sequence_is_correct():
    from test_sign_recognition import BONJOUR_SEQUENCE
    assert BONJOUR_SEQUENCE == ["B", "O", "N", "J", "O", "U", "R"]

def test_zoom_reset():
    z = ZoomManager()
    z.scale = 2.5
    z.reset()
    assert z.scale == 1.0
    assert z.start_distance is None
\end{lstlisting}

\newpage

% ============================================================
% 5. UTILISATION DE L'IA
% ============================================================
\section{Utilisation de l'intelligence artificielle}
\label{sec:ia}

\begin{iabox}
L'IA utilis\'ee est \textbf{Claude Code} (Anthropic, mod\`ele \textit{claude-sonnet-4-6}).
Tout code g\'en\'er\'e ou assist\'e par IA est encadr\'e dans le source~:
\lstinline|# ############### CODE IA (Claude) ################|
\end{iabox}

\vspace{0.4cm}

\begin{table}[H]
\centering
\caption{D\'etail des usages de l'IA dans le projet.}
\label{tab:ia}
\small
\begin{tabularx}{\textwidth}{p{4cm} p{5cm} X}
\toprule
\textbf{Ce qui a \'et\'e fait avec l'IA} & \textbf{Prompt / demande} & \textbf{Ce que j'ai compris / modifi\'e} \\
\midrule
Classe abstraite \texttt{GestureAction} &
``Cr\'ee une classe de base ABC pour des actions gestuelles avec \texttt{process()} abstraite
et \texttt{render\_hud()} commune'' &
V\'erifi\'e la compatibilit\'e avec les sous-classes, adapt\'e les signatures \\
\midrule
Refactorisation h\'eritage &
``Fais h\'eriter \texttt{VolumeController}, \texttt{ScreenshotTaker}, \texttt{ZoomManager}
de \texttt{GestureAction}'' &
Valid\'e que la logique m\'etier (pincement, capture, zoom) reste inchang\'ee \\
\midrule
\texttt{config\_loader.py} &
``Cr\'ee un loader JSON avec deepcopy et fallback sur valeurs par d\'efaut'' &
Identifi\'e et corrig\'e le bug de \textit{shallow copy}~: \texttt{.copy()} ne suffit pas
pour les dicts imbriqu\'es \\
\midrule
29 tests unitaires &
``G\'en\`ere des tests sans webcam pour \texttt{fingers\_up}, \texttt{recognize\_sign},
\texttt{ZoomManager}, config'' &
Ajout\'e les tests sur la s\'equence BONJOUR et les patterns it\'eratifs \\
\midrule
Documentation LaTeX &
``G\'en\`ere le fichier \texttt{.tex} de documentation dans le style du template'' &
R\'edig\'e tout le contenu (contexte, analyse, limites), adapt\'e les figures TikZ \\
\midrule
Correction bug \texttt{HandDetector} &
(d\'etection d'erreur) &
Compris pourquoi \texttt{\_\_init\_\_} ne passait pas les param\`etres \`a MediaPipe \\
\bottomrule
\end{tabularx}
\end{table}

\vspace{0.3cm}
\noindent\textbf{Taux d'utilisation estim\'e~:} $\approx$\,40\,\% (architecture POO, tests, documentation).
La logique de d\'etection gestuelle, la boucle de modes, l'int\'egration MediaPipe/OpenCV
et les patterns de signes (dont la correction it\'erative) ont \'et\'e \'ecrits et ma\^itris\'es
par l'auteur.

\newpage

% ============================================================
% 6. INSTALLATION ET LANCEMENT
% ============================================================
\section{Installation et lancement}
\label{sec:install}

\subsection{Pr\'erequis}

\begin{itemize}
  \item Python 3.11+ (test\'e sur macOS 14 Sonoma)
  \item Webcam fonctionnelle
  \item Permissions macOS~: \textit{Accessibilit\'e} (PyAutoGUI) et
        \textit{Enregistrement d'\'ecran} (MSS) dans
        \texttt{R\'eglages syst\`eme > Confidentialit\'e et s\'ecurit\'e}
\end{itemize}

\subsection{Installation}

\begin{lstlisting}[language=bash, caption={Installation des d\'ependances.}]
# Depuis le dossier HandControlPC--version/
source ../.venv/bin/activate
pip install -r requirements.txt
\end{lstlisting}

\subsection{Lancement}

\begin{lstlisting}[language=bash, caption={Commandes de lancement.}]
# Controle gestuel principal (volume, screenshot, zoom)
python main.py

# Reconnaissance de signes / sequence BONJOUR
python test_sign_recognition.py

# Tests unitaires
python -m pytest tests/ -v
\end{lstlisting}

\subsection{Touches clavier}

\begin{table}[H]
\centering
\begin{tabular}{ll}
\toprule
\textbf{Touche} & \textbf{Action} \\
\midrule
\texttt{ESC} & Quitter le programme \\
\texttt{R} & R\'einitialiser la s\'equence BONJOUR \\
\bottomrule
\end{tabular}
\end{table}

\newpage

% ============================================================
% 7. LIMITES ET PERSPECTIVES
% ============================================================
\section{Limites et perspectives}
\label{sec:limites}

\subsection{Limites actuelles}

\begin{itemize}
  \item \textbf{macOS uniquement}~: \texttt{osascript} (volume) et \texttt{afplay} (son)
        sont sp\'ecifiques \`a macOS. Un portage Windows/Linux n\'ecessiterait des
        alternatives (\texttt{pycaw}, \texttt{pactl}).
  \item \textbf{Luminosit\'e}~: les performances de MediaPipe se d\'egradent en faible
        \'eclairage ou avec des arri\`ere-plans complexes.
  \item \textbf{Signes statiques uniquement}~: les gestes dynamiques (J en LSF est normalement
        un mouvement) ne sont pas support\'es --- on utilise une approximation statique.
  \item \textbf{Mono-utilisateur}~: le syst\`eme ne g\`ere qu'une seule s\'equence BONJOUR
        \`a la fois, sans profil utili\-sateur.
\end{itemize}

\subsection{Pistes d'am\'elioration}

\begin{itemize}
  \item Ajout du d\'eplacement de la souris par geste (index pointant vers l'\'ecran).
  \item Support Windows/Linux via des APIs cross-platform.
  \item Mode apprentissage~: enregistrer de nouveaux patterns de gestes sans modifier le code.
  \item Interface de configuration graphique (tkinter).
  \item Syst\`eme de logs persistants dans un fichier (module \texttt{logging}).
  \item D\'etection de gestes dynamiques avec une fen\^etre temporelle sur les landmarks.
\end{itemize}

\newpage

% ============================================================
% REFERENCES
% ============================================================
\section*{R\'ef\'erences}
\addcontentsline{toc}{section}{R\'ef\'erences}

\begin{description}
  \item[MediaPipe (2023)] Google. \textit{MediaPipe Hands~: On-device Real-time Hand Tracking}.\\
    \url{https://mediapipe.readthedocs.io/en/latest/solutions/hands.html}

  \item[OpenCV (2024)] OpenCV Team. \textit{OpenCV --- Open Source Computer Vision Library}.\\
    \url{https://opencv.org/}

  \item[Python ABC] Python Software Foundation.
    \textit{Module \texttt{abc} --- Abstract Base Classes}.\\
    \url{https://docs.python.org/3/library/abc.html}

  \item[pytest (2024)] \textit{pytest --- helps you write better programs}.\\
    \url{https://docs.pytest.org/}

  \item[MSS (2024)] \textit{MSS --- An ultra fast cross-platform multiple screenshots module}.\\
    \url{https://python-mss.readthedocs.io/}
\end{description}

\end{document}
